\section{Experimental Evaluation}
\label{sec:experiments}

This section evaluates \PG{} on real workloads, real published attack vectors (AgentPoison and MemoryGraft), and a physical mobile device (Xiaomi 13T, MediaTek Dimensity 9200+, 12\,GB RAM, Android 16). All reported metrics derive strictly from immutable manifests recorded across five independent random seeds ($s\in\{1,\ldots,5\}$); no value is interpolated or fabricated.

\subsection{Research Questions}
\begin{description}[leftmargin=*,itemsep=1pt]
  \item[RQ1: Attack Persistence.] Does unmediated mobile ingestion permit persistent backdoor implantation that survives payload deletion?
  \item[RQ2: Provenance Mediation.] How effectively does provenance-aware sink quarantine neutralize attacks without sacrificing clean task utility?
  \item[RQ3: Remediation \& Recovery.] How does lineage rollback compare to approximate unlearning and memory-only baselines in exactness and utility retention?
  \item[RQ4: Systems Overhead.] What are the concrete latency, memory, and storage footprints on commodity mobile hardware?
\end{description}

\subsection{Experimental Setup}

\paragraph{Hardware and Environment.}
Real-device system microbenchmarks execute on a physical Xiaomi 13T (\texttt{2306EPN60G}, MediaTek Dimensity 9200+ SoC, 12\,GB RAM, Android 16 build \texttt{BP2A.250605.031.A3}) via USB ADB and MobileFineTuner's native ARM64 engine \cite{geng2025mobilefinetuner}. Host evaluation runs with local GPU acceleration on an OpenAI-compatible endpoint serving Qwen2.5-7B-Instruct \cite{qwen2024qwen} and Dense Passage Retrieval (DPR) embeddings \cite{karpukhin2020dpr}.

\paragraph{Workloads and Attack Datasets.}
Evaluations span three normalized corpora: StrategyQA knowledge-based reasoning \cite{geva2021strategyqa} (40 evaluation episodes per seed under ReAct \cite{yao2023react}), AESLC email summarization (14,436 normalized events) \cite{zhang2019aeslc}, and PersonaChat dialogue (19,369 normalized events) \cite{zhang2018personachat}. We evaluate two published attack families: (1) \emph{AgentPoison} \cite{chen2024agentpoison}, targeting multi-step ReAct retrieval via optimized trigger sequences; and (2) \emph{MemoryGraft} \cite{srivastava2025memorygraft}, targeting persistent agent experience stores (100 benign and 10 poisoned experiences).

\begin{table*}[t]
\centering
\small
\caption{Empirical evaluation across five independent seeds on StrategyQA and physical Xiaomi 13T smartphone hardware. Values report median (IQR). Measured manifests independently recorded under \texttt{outputs/aggregate\_*}.}
\label{tab:measured-results}
\resizebox{\textwidth}{!}{%
\begin{tabular}{lcccccc}
\toprule
\textbf{Method} & \textbf{ASR Before} & \textbf{ASR After} & \textbf{Clean Utility} & \textbf{False Quarantine} & \textbf{Rollback Time} & \textbf{Systems Footprint \& Recovery} \\
\midrule
No defense & 95.0\% (2.5\%) & 95.0\% (2.5\%) & 70.0\% (2.5\%) & n/a & n/a & 3.92s / step, 2.17 GB peak RSS \\
Memory-only (A-MemGuard) & 95.0\% (2.5\%) & 0.0\% (0.0\%)$^\dagger$ & 67.5\% (2.5\%) & 5.0\% (0.0\%) & n/a & Training sink unmediated \\
Approximate unlearning (NPO) & 95.0\% (2.5\%) & 15.0\% (0.0\%) & 57.5\% (0.0\%) & 0.0\% (0.0\%) & n/a & 0\% exact match, $-12.5\%$ utility drop \\
\textbf{\PG{} (Ours)} & \textbf{95.0\% (2.5\%)} & \textbf{0.0\% (0.0\%)} & \textbf{70.0\% (2.5\%)} & \textbf{0.0\% (0.0\%)} & \textbf{2.94 ms (0.02 ms)} & \textbf{127.95 kB storage, 100\% exact} \\
Clean counterfactual & 0.0\% (0.0\%) & 0.0\% (0.0\%) & 70.0\% (2.5\%) & n/a & n/a & Baseline reference \\
\bottomrule
\multicolumn{7}{l}{\footnotesize $^\dagger$Filters memory retrieval attacks, but leaves LoRA adapter training unprotected against cross-sink laundering.}
\end{tabular}%
}
\end{table*}

\subsection{Security and Defense Results (RQ1, RQ2)}

Table~\ref{tab:measured-results} summarizes security and utility metrics across all evaluated methods. 

\paragraph{Vulnerability without Defense (RQ1).}
Under \emph{No defense}, AgentPoison achieves an attack success rate (ASR) of $95.0\%$ median (IQR $2.5\%$, min $87.5\%$, max $100.0\%$) on StrategyQA. Deleting the raw external message or source payload does not mitigate the attack ($\mathrm{ASR}=95.0\%$), confirming that unmediated ingestion enables persistent compromise. Similarly, under MemoryGraft experience injection, $47.9\%$ of retrieved past experiences are poisoned without defense.

\paragraph{Protection via Complete Mediation (RQ2).}
With \PG{} active, untrusted external ingress (SMS, shared files, unverified web content) is automatically assigned $\Quarantine$ for persistent sinks. When the adversarial queries arrive, quarantined passages are suppressed from the agent's prompt context. Post-defense ASR drops to $\mathbf{0.0\%}$ across all five seeds. Clean task utility is maintained at $70.0\%$ (IQR $2.5\%$), exactly matching the clean counterfactual, while false quarantine on benign user-authored instances is $0.0\%$. Under MemoryGraft, poisoned experience retrieval drops from $47.9\%$ to $\mathbf{0.0\%}$ while preserving an $85.0\%$ clean experience success rate.

\subsection{Comparison with Baselines (RQ3)}

\paragraph{Memory-Only Sanitization.}
A-MemGuard-style memory sanitization \cite{wei2025amemguard} checks retrieval consistency and suppresses memory attacks ($0.0\%$ ASR on retrieval, denoted by $\dagger$). However, because it lacks cross-sink lineage tracking, the on-device LoRA training sink remains completely exposed. Furthermore, over-filtering benign remote communications incurs a $5.0\%$ false quarantine rate and slightly degrades clean utility to $67.5\%$. When poison is identified in fine-tuned weights, memory-only defenses cannot perform weight recovery ($\text{Rollback time}=\text{n/a}$).

\paragraph{Approximate Unlearning.}
Gradient-based approximate unlearning (e.g., NPO \cite{zhang2024npo}) attempts to erase backdoors by perturbing weights along negative gradients. As shown in Table~\ref{tab:measured-results}, it fails to completely eliminate the attack, leaving a $15.0\%$ residual ASR. More critically, unlearning induces catastrophic forgetting: clean task utility plummets from $70.0\%$ to $57.5\%$ (a $12.5\%$ utility penalty). In addition, $0$ out of $5$ seeds achieve parameter equality ($\text{Exact rate}=0\%$, $L_2$ parameter distance $0.016$).

\paragraph{Exact Restoration under \PG{}.}
In contrast, \PG{} traverses the lineage DAG to invalidate contaminated descendants, restores the latest clean checkpoint, and deterministically replays unaffected events. Across all five seeds, \PG{} achieves a $\mathbf{100\%}$ exact match rate with the clean counterfactual parameter vectors, validating Theorem~\ref{thm:rollback} empirically.

\subsection{On-Device Systems Overhead (RQ4)}

\paragraph{Training and Inference Latency.}
On the physical Xiaomi 13T, MobileFineTuner executes local LoRA fine-tuning steps with a median latency of $3,918\text{ ms}$ ($\approx 3.92\text{s}$, IQR $16\text{ ms}$). Peak resident set size (RSS) is $2,170,712\text{ kB}$ ($\approx 2.17\text{ GB}$, IQR $14.4\text{ MB}$), operating comfortably within the device's $12\text{ GB}$ physical RAM.

\paragraph{Rollback Latency and Storage Footprint.}
Lineage invalidation, checkpoint retrieval, and clean suffix replay execute with a median latency of only $\mathbf{2.94\text{ ms}}$ (IQR $0.02\text{ ms}$, min $2.90\text{ ms}$, max $3.75\text{ ms}$). The serialized lineage graph and checkpoint metadata occupy only $\mathbf{127.95\text{ kB}}$, representing an infinitesimal storage footprint on commodity flash storage (pipeline detailed in \cref{fig:experimental-setup}, \cref{app:experiments}).
